% === ABSTRACT ===
\begin{abstract}
Population aging is accelerating globally: the proportion of adults aged 60+ is projected to double by 2050, while physiotherapist-to-population ratios remain critically low across low- and middle-income countries. We present \system{}, a multimodal AI rehabilitation system designed for elderly users, and report a \emph{systems feasibility} study. The system integrates four components: (1)~real-time whole-body 3D pose estimation via RTMW3D-L~\cite{ref_rtmw3d} from MMPose~\cite{ref_mmpose}, producing 133 keypoints (body, hands, face in COCO-WholeBody format) with quaternion-based kinematics for multi-plane joint angle computation; (2)~multi-indicator facial state detection through OpenFace~3.0 Action Unit analysis, PSPI-inspired pain estimation~\cite{ref_pspi}, PERCLOS-based fatigue monitoring~\cite{ref_perclos}, and engagement detection~\cite{ref_engagement}; (3)~voice-interactive coaching powered by Large Language Models with Retrieval-Augmented Generation, grounding exercise guidance in a clinical knowledge base; and (4)~a six-dimension scoring framework incorporating SPARC smoothness~\cite{ref_sparc}, constrained DTW alignment, and temporal compensation detection. The complete pipeline---pose estimation, facial analysis, scoring, and coaching---sustains 129.7~FPS on a single consumer GPU, establishing that real-time multimodal rehabilitation feedback is computationally feasible without specialized depth sensors. To our knowledge, \system{} is the first integrated rehabilitation system combining whole-body 3D pose estimation, multi-indicator facial state detection, voice-interactive LLM coaching, and clinically-inspired kinematic scoring in a single real-time framework. We identify clinical validation and scoring calibration as the necessary next steps.
\end{abstract}

\begin{IEEEkeywords}
rehabilitation, pose estimation, elderly care, multimodal AI, LLM coaching, pain detection, whole-body tracking, RTMW3D, SPARC, facial action units
\end{IEEEkeywords}
